\section{Introduction}
\label{sec:introduction}

Large language models (LLMs) are increasingly deployed as autonomous
agents capable of interacting with external tools and
APIs~\citep{schick2024toolformer, yao2023react}. Rather than generating
text in isolation, these agents retrieve documents, query databases,
send messages, and trigger downstream workflow actions with
real-world side effects. A growing body of work surveys this
paradigm~\citep{wang2024agent_survey}, and several orchestration
frameworks have emerged to structure agent behavior as explicit workflow
graphs: LangGraph~\citep{langgraph_docs},
CrewAI~\citep{crewai_docs}, AutoGen~\citep{wu2023autogen}, and Google
ADK~\citep{google_adk_docs} each represent computation steps as nodes
and allowable transitions as edges.

Despite this rich structural representation, safety enforcement in
current agent systems remains predominantly a \emph{runtime} concern.
Tools such as NeMo Guardrails~\citep{nemo_guardrails} and
LlamaGuard~\citep{llamaguard} intercept individual LLM calls or tool
invocations and apply content-level or policy-level filters at execution
time. While valuable for catching toxic outputs and prompt injection
attempts, these runtime approaches have three fundamental limitations:
(i)~they impose per-call latency overhead, (ii)~they detect violations
only when the offending execution path is actually exercised, and
(iii)~they provide no coverage guarantees over paths that remain
untriggered during testing.

\paragraph{Motivating example.}
Consider an email triage agent with the following workflow: incoming
emails are classified by intent, routed to either an urgent or normal
handler, drafted, and sent. During iterative development, a developer
adds a \texttt{draft\_response} node on the normal-priority path but
forgets to connect it to the \texttt{send} node. The result is a
\emph{dead end}: normal-priority emails enter the draft stage and
silently halt. A runtime guardrail would not catch this defect unless
the normal-priority path is exercised with an email that triggers it.
A static analysis of the workflow graph, however, immediately flags the
dead-end node and produces a witness trace:
\texttt{\_\_start\_\_ $\to$ classify $\to$ router $\to$
normal\_handler $\to$ draft\_response} (stuck).

\paragraph{Research gap.}
This class of defect---a topological error in the workflow
graph cannot be addressed by content-level runtime guardrails. In
principle, classical model checking techniques could detect such
errors: reachability analysis, deadlock detection, and temporal logic
verification are well-established in the formal methods
literature~\citep{clarke1999modelchecking, baier2008modelchecking}.
However, general-purpose model checkers such as
SPIN~\citep{holzmann1997spin} and NuSMV~\citep{cimatti2002nusmv}
require manual translation of the system under analysis into a
dedicated modeling language---a prohibitive overhead for agent
developers iterating on workflow designs. Similarly, the business
process management (BPM) community has a long tradition of workflow
soundness verification~\citep{vanderaalst2011bpm}, but these methods
target visual modeling tools and standardized notations (e.g., BPMN,
Petri nets), not the programmatic API objects through which agent
frameworks define their graphs. To date, no system automatically
extracts and statically verifies agent workflow graphs directly from
framework source code.

\paragraph{Proposed approach.}
This paper argues for \emph{pre-deployment static verification} of
agent workflow graphs. The key observation is that modern agent
frameworks expose workflow structure through their APIs, making the
graph amenable to automated extraction and analysis. When agent behavior
is graph-structured, safety properties reduce to reachability,
isolation, and temporal constraints over the topology properties that
can be checked exhaustively without running the
agent~\citep{clarke1999modelchecking}.

This paper presents Agentproof, a practical system for static analysis
of agent workflows that focuses on properties (a)~expressible over the
workflow graph and (b)~independent of LLM text generation semantics.
Unlike general-purpose model checkers, Agentproof extracts the analysis
model automatically from framework source code, requiring no manual
modeling effort.

\paragraph{Contributions.}
\begin{itemize}
  \item \textbf{The first real-world base-rate study} of structural
  (topology) defects in agent workflows: 930 workflows mined from 253
  public GitHub repositories across all four frameworks (LangGraph,
  AutoGen, CrewAI, ADK), with a 120-workflow subsample validated against
  independent LLM-reconstructed ground-truth graphs and adversarial defect
  triage (Section~\ref{sec:realworld}). The finding is negative, and
  because nobody had measured it the rarity is itself the contribution:
  the validated topology-defect rate is ${\approx}0\%$ (2 of 187 flags
  genuine; 43\% extraction artifacts, 50\% intentional design), and real
  workflows are predominantly small static linear pipelines (median two
  non-sentinel nodes).
  \item \textbf{\textsc{Agentproof}}, the measurement instrument that made
  the study possible: automatic extraction of a unified abstract graph
  model from four heterogeneous agent frameworks with no hand-written
  formal model (unlike SPIN/NuSMV, which require Promela/SMV), six
  structural checks with witness-trace generation, and a seven-form
  temporal policy DSL over the safety fragment of LTL, compiled to DFAs
  and evaluated both statically and at
  runtime~(Sections~\ref{sec:graph_model}--\ref{sec:verification_methods}).
  \item \textbf{A quantified diagnosis} with an actionable failure-mode
  taxonomy: extraction fidelity, not the check algorithms, is the binding
  constraint on static-analysis utility. On real code node
  precision/recall stays high (${\sim}0.91/0.85$) but edge recall drops
  to 0.64 and node-kind accuracy to 0.83 (both near-perfect on stubs),
  and fidelity is strongly framework-dependent---from 0.77 edge recall on
  AutoGen down to 0.37 on ADK, whose topology is most implicit. Every
  structural flag we triaged is an extraction artifact, and such flags
  arise only on LangGraph (78\% of its workflows) because the other
  frameworks' extractors impose a connected topology by construction.
  \item \textbf{A defect-independent job} for the static machinery: the
  graph\,$\times$\,DFA product proves which temporal policies can never
  reach a violation on a given workflow, so their runtime monitors can be
  provably pruned at enforcement time---positioning static verification
  as a monitor-selection layer that bridges to leaner runtime shields.
\end{itemize}

\paragraph{Scope.}
The proposed approach validates workflow design decisions early: ``can this
graph reach a destructive tool?'' or ``must human review occur between
two sensitive actions?'' It does not attempt to prove semantic
properties of LLM outputs (e.g., truthfulness), which remain
inherently non-deterministic. This boundary is discussed precisely in
the threat model (Section~\ref{sec:threat_model}).
